% chktex-file 8
% chktex-file 12
% chktex-file 18
\beginsong{Marsyas a Apollon}[by={Marsyas}]

\beginverse{}
\ch{A}{}{}{}Ta krásná dívka, \ch{Bmi}{}{}{}co se bojí o \ch{E}{}{}{}svoji krásu
\ch{A}{}{}{}Athéna \ch{F\#mi}{}{}{}jméno má \ch{D}{}{}{}za starých \ch{A}{}{}{}dávných časů
\ch{A}{}{}{}Odhodí flétnu, \ch{Bmi}{}{}{}hrát nejde s \ch{E}{}{}{}nehybnou tváří
\ch{A}{}{}{}Ten, kdo ji \ch{F\#mi}{}{}{}najde dřív, \ch{D}{}{}{}tomu se \ch{A}{}{}{}přání zmaří
\endverse{}

\beginchorus{}
\ch{F\#mi}{}{}{}Tak si Marsyas \ch{E}{}{}{}mámen flétnou vě\ch{A}{}{}{}ří, že \ch{D}{}{}{}musí přetnout
\lrep{} \ch{A}{}{}{}jedno pravidlo, \ch{E}{}{}{}sázku a \ch{D}{}{}{}hrát\ch{F\#mi}{}{}{} líp \ch{Bmi}{}{}{}než bů\ch{A}{}{}{}h \rrep{}
\endchorus{}

\beginverse{}
Bláznivý nápad, snad nejvýš Marsyas míří
Apollón souhlasí, oba se s trestem smíří
Král Midas má říct, kdo je lepši, Apollón zpívá
O život soupeří, jen jeden vítěz bývá
\endverse{}

\beginchorus{}
Tak si Marsyas mámen flétnou věří a musí přetnout
\lrep{} jedno pravidlo, sázku a hrát líp než bůh \rrep{}
\endchorus{}

\beginverse{}
Obrátí nástroj, už ví, že nebude chválen
Prohrál a zápolí podveden vůlí krále
Sám v tichém hloučku, sám na strom připraví ráhno
Satyra k hrůze všech zaživa z kůže stáhnou
\endverse{}

\beginchorus{}
Tak si Apollón změřil síly, každý se musel mýlit
\lrep{} Nikdo nemůže kouzlit a hrát líp než bůh \rrep{}
\endchorus{}

\beginverse{}
Ta krásná dívka, co se bála o svoji krásu
Dárkyně moudrosti za starých dávných časů
Teď v tichém hloučku, v jejích rukou úroda, spása
Athéna jméno má, chybí jí tvář a krása
\endverse{}

\beginchorus{}
Dy dy dy \ldots{}..
\endchorus{}

\endsong{}
